\documentclass[11pt]{article}
\usepackage{coursenotes}
\begin{document}
\handoutheader{H3}{Staggered Adoption}{Cohorts, clean comparisons, and what two-way fixed effects really averages}{Meeting 10 (Panel data and DiD)}

\begin{decision}
Real rollouts are staggered: some stores adopt in March, some in May, some never. The two-way fixed-effects regression that
works for one adoption date also compares late adopters with early ones that are already treated, without saying so. If effects grow or fade
over time, its answer can be far from any average effect, even of the wrong sign. This handout builds the estimate from clean
comparisons instead.
\end{decision}

\begin{goals}
  \item define group-time effects with adoption-date potential outcomes;
  \item identify each group-time effect with never-treated or not-yet-treated comparisons, and handle anticipation;
  \item show why already-treated units are contaminated controls, and what TWFE averages;
  \item aggregate group-time effects by event time and overall, and build event studies from clean comparisons;
  \item describe the design-based alternative: randomised adoption dates.
\end{goals}

%=====================================================================================
\sectionone

\subsection{Setup}

Unit $i$ first receives treatment in period $G_i \in \{2, \dots, T\} \cup \{\infty\}$ and stays treated (no reversal); $D_{it} = \ind{t \ge G_i}$. Potential
outcomes are indexed by adoption date: $Y_{it}(g)$, with $Y_{it}(\infty)$ if never treated; we observe $Y_{it} = Y_{it}(G_i)$. Units sharing an adoption
date form a \emph{cohort}. The building block is the \emph{group-time average treatment effect on the treated}
\[
  \ATT(g, t) = \E[Y_{it}(g) - Y_{it}(\infty) \mid G_i = g], \qquad t \ge g .
\]
Every reported number (event-study coefficients, overall averages) is an average of these.

\begin{assum}{No anticipation}{noant}
$Y_{it}(g) = Y_{it}(\infty)$ for all $t < g$.
\end{assum}
\begin{assum}{Parallel trends with the never-treated}{ptnev}
For every cohort $g$ and $t \ge g$: $\E[Y_{it}(\infty) - Y_{i,g-1}(\infty) \mid G_i = g] = \E[Y_{it}(\infty) - Y_{i,g-1}(\infty) \mid G_i = \infty]$.
\end{assum}
\begin{assum}{Parallel trends with the not-yet-treated}{ptnyt}
The same equality holds with $G_i = \infty$ replaced by $G_i = s$ for every $s > t$.
\end{assum}

\subsection{Identification from clean comparisons}

\begin{prop}{Group-time effects}{gt}
Under no anticipation and parallel trends with the never-treated, for $t \ge g$,
\[
  \ATT(g, t) = \{\E[Y_{it} \mid G=g] - \E[Y_{i,g-1} \mid G=g]\} - \{\E[Y_{it} \mid G=\infty] - \E[Y_{i,g-1} \mid G=\infty]\}.
\]
Under parallel trends with the not-yet-treated, the same holds with the never-treated replaced by $\{i : G_i > t\}$.
\end{prop}
\begin{proof}
For cohort $g$ the first brace is $\E[Y_{it}(g) - Y_{i,g-1}(\infty) \mid G=g]$ by no anticipation, which is $\ATT(g,t)$ plus cohort $g$'s untreated change; the
second brace is the comparison group's untreated change (not-yet-treated units are untreated at $t$ and $g - 1$ by no anticipation); parallel trends
equates the two untreated changes.
\end{proof}

Each $\ATT(g,t)$ is a two-by-two DiD with a \emph{clean} control, whose treatment status does not change between the base period and $t$, and with the
base fixed at $g - 1$. This is the estimator of Callaway and Sant'Anna (2021).

\paragraph{Anticipation moves the base.} If units change behaviour $k$ periods before adoption (training staff off the floor), the base must be $g - k - 1$,
and a later cohort stops being a clean control once \emph{its} anticipation begins. Build the table of eligible controls and base periods from the calendar
before computing any effect.

\subsection{Why already-treated units are contaminated controls}

\begin{prop}{Contamination}{contam}
Let cohort $k$ adopt before cohort $l$, and use $k$ as the control for $l$ between periods $t_0 < l \le t_1$ with $t_0 \ge k$. Under parallel trends for both,
\[
  \text{DiD}_{l\ \text{vs}\ k}(t_0, t_1) = \ATT(l, t_1) - [\ATT(k, t_1) - \ATT(k, t_0)].
\]
\end{prop}
\begin{proof}
Cohort $l$'s change is $\ATT(l, t_1)$ plus its untreated change (no anticipation at $t_0$). Cohort $k$ is treated in both periods, so its change is
$\ATT(k,t_1) - \ATT(k,t_0)$ plus its untreated change. Parallel trends equates the untreated changes; subtract.
\end{proof}

The bias is the \emph{change} in the control cohort's own effect: downward if effects grow, upward if they fade, zero if constant.

\subsection{What two-way fixed effects averages}

\begin{prop}{The Goodman-Bacon decomposition (statement)}{bacon}
With staggered adoption, the TWFE coefficient on $D_{it}$ equals a weighted average of all two-by-two DiDs in the panel: each cohort against the
never-treated, earlier against later cohorts before the later ones adopt, and later against earlier cohorts after the earlier ones adopt. The weights are
non-negative, sum to one, and grow with group sizes and with the variance of treatment within each pair; they are set by the adoption calendar, not by
any estimand.
\end{prop}
\begin{proof}
A Frisch--Waugh--Lovell decomposition of the variance of the two-way demeaned $D_{it}$; see Goodman-Bacon (2021) for the weights.
\end{proof}

The ``later versus earlier'' terms are the contaminated comparisons of Result~\ref{prop:contam}. In terms of the underlying $\ATT(g,t)$, some implicit weights
can be negative, so TWFE can lie outside the range of every group-time effect. With one adoption date there are no such terms, which is why TWFE is fine in
that case.

\subsection{Aggregation and event studies}

\paragraph{By event time.} For $e \ge 0$, $\ATT(e) = \sum_g \Pr(G = g \mid g + e \le T)\,\ATT(g, g+e)$: the average effect $e$ periods after adoption, across
cohorts observed that long. This is what a rollout decision needs: a unit adopting tomorrow starts at $e = 0$.

\paragraph{Overall.} Average all treated $\ATT(g,t)$ cells weighted by cohort size, or average within cohort first. It depends on how long each cohort has been
observed.

\paragraph{Event studies.} Pre-period placebo contrasts (cohort $g$'s one-period change relative to the control, before $g$) should be near zero. Build event
studies from clean comparisons: event-time dummies in a TWFE regression inherit the contaminated comparisons and can show spurious pre-trends.

\paragraph{Imputation.} An alternative estimator fits unit and period effects on untreated observations only (never-treated units and not-yet-treated periods),
imputes $Y_{it}(\infty)$ for treated cells, and averages $Y_{it} - \hat Y_{it}(\infty)$ with chosen weights (Borusyak, Jaravel and Spiess, 2024). It uses all
pre-periods efficiently and never uses treated cells as controls.

\subsection{Inference}

Group-time estimates share control units and base periods, and a unit's observations are correlated over time. Cluster by unit (or by the level of adoption);
for simultaneous bands across event times, use the multiplier bootstrap over units. The number of clusters is the number of units, not unit-periods.

\subsection{The design-based alternative: randomise the dates}

\begin{prop}{Randomised adoption dates}{random}
If the adoption date of each unit is assigned at random from a fixed set of dates (independently of potential outcomes), then for each $t$ and each pair of
cohorts $g \le t < s$, the difference in mean outcomes at $t$ between cohort $g$ and cohort $s$ is unbiased for $\E[Y_{it}(g) - Y_{it}(s)]$, and under no
anticipation, for $\E[Y_{it}(g) - Y_{it}(\infty)]$ over the randomised population.
\end{prop}
\begin{proof}
Random assignment of dates makes cohort membership independent of the potential-outcome vector, so each cohort's mean at $t$ is unbiased for the
population mean of $Y_{it}(g)$ (Meeting~1). No anticipation gives $Y_{it}(s) = Y_{it}(\infty)$ for $t < s$.
\end{proof}

No parallel-trends assumption is needed: the comparison is identified by design, as in an experiment, and the design-based inference of Meeting~2
applies with the randomised dates (Athey and Imbens, 2022). ``Stagger the rest at random'' is therefore the cheapest upgrade to any rollout plan.

\begin{pitfall}[staggered DiD errors]
TWFE with staggered adoption; event-time dummies in TWFE; a base period inside an anticipation window; later cohorts used as controls after their
anticipation begins; choosing which units adopt first by expected benefit when the order could have been randomised.
\end{pitfall}

%=====================================================================================
\sectiontwo

\paragraph{Setting.} FitLife extends app booking in waves. Club group E adopts in month 2, group L in month 3, group N not in the window. The outcome is
classes booked per club per day. A common trend of $+1$ per month runs through all groups; the true effect is $6$ in the first month after adoption and
$3$ in the second (members try the app, then some revert).
\begin{center}
\begin{tabular}{lccc}
\toprule
& Month 1 & Month 2 & Month 3 \\
\midrule
E (adopts month 2) & 40 & \textbf{47} & \textbf{45} \\
L (adopts month 3) & 50 & 51 & \textbf{58} \\
N (never) & 60 & 61 & 62 \\
\bottomrule
\end{tabular}
\end{center}

\paragraph{Clean comparisons.} $\ATT(E,2) = (47 - 40) - (61 - 60) = 6$; $\ATT(E,3) = (45 - 40) - (62 - 60) = 3$; $\ATT(L,3) = (58 - 51) - (62 - 61) = 6$. Event-time path
$(\ATT(0), \ATT(1)) = (6, 3)$; overall average $5$.

\paragraph{The contaminated comparison.} L against E, months $2 \to 3$: $(58 - 51) - (45 - 47) = 9$. Result~\ref{prop:contam} predicts
$\ATT(L,3) - [\ATT(E,3) - \ATT(E,2)] = 6 - (3 - 6) = 9$: E's effect \emph{fell} by $3$, and the comparison credited the fall to L as extra gain.

\paragraph{TWFE.} With equal group sizes the Goodman-Bacon weights are $\tfrac13$ (E vs N), $\tfrac13$ (L vs N), $\tfrac16$ (E vs L before L adopts) and $\tfrac16$ (L vs E after
E adopts). The four two-by-two estimates are $4.5$, $6$, $6$ and $9$, so
$\hat\beta_{\text{TWFE}} = \tfrac13(4.5) + \tfrac13(6) + \tfrac16(6) + \tfrac16(9) = 6$, which least squares on the nine cells confirms. TWFE overstates the overall average of
$5$ and reports the first-month effect as if it lasted.

\paragraph{The decision.} If the app is worth extending only when it adds at least $5.5$ classes a day in steady state, TWFE says yes and the event-time path
says the gain fades to $3$: no, unless the fade can be reversed.

\begin{exercise}[computation]
Add a fourth month in which both E's and L's effects are $3$ (E: 46, L: 56, N: 63). (a) Compute $\ATT(L,4)$ with N as the control and base month 2.
(b) Compute the same comparison with E as the control. (c) Explain the difference with Result~\ref{prop:contam}.
\end{exercise}
\answer{(a) $(56 - 51) - (63 - 61) = 3$. (b) $(56 - 51) - (46 - 47) = 5 + 1 = 6$. (c) E is treated in both month 2 and month 4, and its own effect fell from $6$ to
$3$ between them; Result~\ref{prop:contam} gives $\ATT(L,4) - [\ATT(E,4) - \ATT(E,2)] = 3 - (3 - 6) = 6$. The already-treated control credits E's fade to L.}

\begin{exercise}[judgement]
FitLife's operations team wants the next wave to go to the clubs ``most ready for the app''. What does that do to the evaluation, and what would you
propose?
\end{exercise}
\answer{Readiness probably predicts both adoption timing and the size of the effect (and perhaps trends), which threatens parallel trends and makes
early-cohort effects unrepresentative of later ones. Propose randomising the order of the remaining waves (possibly within readiness strata), so cohort
comparisons are identified by design (Result~\ref{prop:random}).}

%=====================================================================================
\sectionthree

\subsection*{Annotated reading}
\begin{readinglist}
\reading{Goodman-Bacon (2021), ``Difference-in-differences with variation in treatment timing'', \emph{Journal of Econometrics}}{The decomposition of
TWFE into two-by-two comparisons}{Sections 1--3}{2}
\reading{Callaway and Sant'Anna (2021), ``Difference-in-differences with multiple time periods'', \emph{Journal of Econometrics}}{Group-time effects,
aggregation and inference}{Sections 1--4}{2}
\reading{Sun and Abraham (2021), ``Estimating dynamic treatment effects in event studies with heterogeneous treatment effects'', \emph{Journal of
Econometrics}}{Why TWFE event studies mislead, and a fix}{Sections 1--3}{3}
\reading{Borusyak, Jaravel and Spiess (2024), ``Revisiting event-study designs: robust and efficient estimation'', \emph{Review of Economic
Studies}}{The imputation estimator}{Sections 1--3}{3}
\reading{de Chaisemartin and D'Haultf\oe uille (2020), ``Two-way fixed effects estimators with heterogeneous treatment effects'', \emph{AER}}{Negative
weights in TWFE}{Sections I--III}{3}
\reading{Athey and Imbens (2022), ``Design-based analysis in difference-in-differences settings with staggered adoption'', \emph{Journal of
Econometrics}}{Randomised adoption dates and design-based inference}{Sections 1--4}{3}
\reading{Roth, Sant'Anna, Bilinski and Poe (2023), ``What's trending in difference-in-differences?'', \emph{Journal of Econometrics}}{A synthesis of
the literature}{All}{2}
\end{readinglist}

\end{document}
